\begin{textsample}{POS Dim 1 – mt – Score 51.00 – mt/mt\_ef\_35.txt}  \label{ex:f1_pos_075}
1.1.1.
Contribuição da Língua Portuguesa para os Anos Finais do Ensino Fundamental O aprendizado de \textbf{uma} língua, especialmente da Língua Portuguesa, oportuniza ao indivíduo o direito de, ao conhecê -la em seus aspectos formais, utilizá -la para promover sua cidadania e tornar -se capaz de \textbf{exigir} direitos e cumprir deveres sabendo como e por \textbf{que} os cumpre ou os \textbf{exige}.
Não há como falar em direitos sem discutir o aprendizado de Línguas.
Não há como falar em cidadão brasileiro sem pensar em como isso está posto em todo e qualquer documento, da certidão de nascimento, à de óbito.
Todos os papéis \textbf{que} assinamos, todos aqueles \textbf{que} cumprimos como “ \textbf{atores} sociais ”, estão vinculados ao saber usar a Língua, mais ou menos formalmente.
Isto \textbf{posto}, vale \textbf{lembrar} \textbf{que} este documento, elaborado em Língua Portuguesa, só terá sentido para quem \textbf{conseguir} lê -lo e compreendê -lo em suas múltiplas linguagens e formas de letramento.
Nosso papel é dar mais \textbf{que} o acesso a norma padrão, é mostrar suas facetas no uso cotidiano e propiciar seu uso nas práticas de linguagem.
Nessa perspectiva, o ensino de Língua Portuguesa amplia e consolida a compreensão das práticas de linguagem, como descritas nas competências gerais do componente curricular, \textbf{que} por sua vez estão atreladas às competências gerais da BNCC.
Não há como pensar em construção de competências básicas, em qualquer área, sem pensar em como se dá sua transmissão \textbf{que}, no Brasil, oficialmente acontece através da Língua Padrão. 1.1.2.
Práticas de Linguagem O ensino da língua Portuguesa tem como princípio e pressuposto a \textbf{abordagem} dos eixos de integração, também considerados pela BNCC, mas \textbf{que} “ devem estar envolvidos em práticas de reflexão ”, \textbf{que} estão organizados e especificados em suas práticas de linguagem ( unidades temáticas ) : a leitura, oralidade, escrita, conhecimentos linguísticos e gramaticais e educação \textbf{literária}, os quais se apresentarão a seguir, \textbf{pontuando} o processo de ensino/aprendizagem de Língua Portuguesa dos anos finais do fundamental : 1.1.2.1.
Leitura Em relação à leitura, a BNCC ( 2017 ) conceitua como práticas de linguagem \textbf{que} decorrem da interação ativa do leitor/ouvinte/espectador com os textos escritos, orais e multissemióticos em diferentes portadores e sua interpretação.
Nesse sentido, a prática da leitura assume \textbf{um} papel fundamental nas aulas de Língua Portuguesa como instrumento para a reflexão e ressignificação do texto no contexto social do estudante.
O tratamento das práticas leitoras na unidade escolar deve considerar o leitor de forma a aceitar seu conhecimento prévio, seja advindo de práticas escolares, seja considerando os textos \textbf{que} circulam em seu meio, ampliando seu repertório com textos \textbf{que} estão presentes em outras esferas sociais, tais como : textos jornalísticos, científicos, \textbf{literários}, artísticos culturais, publicitários, midiáticos, e outros \textbf{que} sirvam de aporte para discussões e formação de opinião acerca de temas como diversidade, política, preconceito e outros temas, pois, como nos ensina Orlandi ( 2001, p. 40 ), > No \textbf{que} \textbf{diz} respeito às diferentes formas de linguagem \textbf{que} constituem o universo simbólico desse aluno, seria \textbf{interessante} \textbf{que}, ao invés de ser \textbf{uma} relação suposta e recusada, ela fosse o ponto de partida, a fonte de hipóteses para estimular e fazer avançar o processo do \textbf{aprendiz}.
A convivência com a música, a pintura, a fotografia, o cinema, com outras formas de utilização do som e com a imagem, assim como a convivência com as linguagens artificiais poderiam nos apontar para \textbf{uma} inserção no universo simbólico \textbf{que} não é a \textbf{que} temos estabelecido na escola [... ] essa articulação \textbf{que} deveria ser explorada no ensino da leitura [... ].
O \textbf{sujeito} no ato da leitura utiliza estratégias para compreensão do texto, \textbf{baseado} em seu conhecimento linguístico e na sua vivência sociocultural, considerando \textbf{que} o leitor não apreende apenas o sentido presente \textbf{que} está no texto, mas também lhe atribui sentidos ( ORLANDI, 2001 ).
Dessa forma, para o desenvolvimento de leituras significativas, o professor deve proporcionar condições para \textbf{que} o estudante entre em contato com textos \textbf{que} lhe permitam o estabelecimento de relações e sentidos, tanto no texto quanto em sua própria vida. 1.1.2.2.
Oralidade Sabe -se \textbf{que}, comumente, algumas escolas brasileiras desconsideram a relevância da modalidade oral para o Ensino da Língua Materna, privilegiando a modalidade escrita.
No entanto, no cenário \textbf{atual} das discussões, reflexões, pesquisas e da BNCC é indispensável o trabalho com a oralidade, considerando sua relevância à formação do \textbf{sujeito} formador de opinião e cidadão de direitos e deveres em \textbf{uma} sociedade letrada.
Reitera -se a importância do trabalho com a oralidade na construção do \textbf{sujeito}.
Em 1998, os Parâmetros Curriculares Nacionais dispunham : > Ensinar língua oral deve significar para a escola possibilitar acesso a usos da linguagem mais formalizados e convencionais, \textbf{que} \textbf{exijam} controle mais consciente e voluntário da enunciação, tendo em vista a importância \textbf{que} o domínio da palavra pública tem no exercício da cidadania ( BRASIL, 1998, p. 65 ).
Essa posição é reiterada pela BNCC, \textbf{que} a coloca como eixo temático em todos os campos, não só da Língua Portuguesa, mas em todas as áreas do conhecimento e desde a educação infantil.
Assim, a escola, responsável por garantir a aprendizagem do estudante, deve oportunizar aos alunos o uso adequado da língua oral nas mais diversas situações de interação verbal.
É importante destacar \textbf{que}, na medida em \textbf{que} se avança nos níveis de estudo, diminui -se o uso pedagógico da oralidade, na contramão do \textbf{que} pedem os documentos citados.
Mesmo fazendo uso informal da língua, a partir do final do 3º ano, os alunos deixam de vivenciar a oralidade em muitas interações sociais.
Outro aspecto \textbf{que} se deve apontar é o fato de \textbf{que} as variações linguísticas, \textbf{que} deveriam ser respeitadas como tais, dão espaço à linguagem formal, sem \textbf{que}, no entanto, seja explicada a necessidade de aprenderem \textbf{uma} nova variação, diferente da \textbf{que} já possuem e atrelada ao uso social da língua como instrumento de ascensão social.
Nesse sentido, é necessário no planejamento pedagógico \textbf{que} os educadores articulem ações de aprendizagem referentes ao uso da língua nas práticas de oralidade, possibilitando ao estudante realizar a comunicação em diversos contextos sociais, sempre estabelecendo a relação entre fala e escrita, pois ambas estão intimamente ligadas e articuladas no processo discursivo.
O trabalho com a oralidade deve ocorrer de modo significativo com atividades \textbf{que} possibilitem a formação integral do estudante, de forma reflexiva sobre as condições de produção dos textos orais \textbf{que} circulam em diferentes mídias e campos de atividade humana.
A produção de textos orais e compreensão dos efeitos de sentido provocados pelos usos de recursos linguísticos e multissemióticos deverão conduzir os alunos aos princípios \textbf{que} regem o funcionamento da língua, considerando o interlocutor e possibilitando o uso adequado da fala em diversas situações de comunicação. 1.1.2.3.
Escrita Neste \textbf{tópico}, \textbf{convém} esclarecer \textbf{que}, assim como os documentos e orientações curriculares elaborados nas últimas décadas, este documento de referência curricular também assume a concepção interacionista da linguagem e a heterogeneidade do \textbf{sujeito} \textbf{historicamente} constituído, buscando, no entanto, incluir as transformações das práticas de linguagem ocorridas neste \textbf{século}, como também as pesquisas desenvolvidas em decorrência da expansão das tecnologias da informação e comunicação.
Isso porque se acredita \textbf{que} o trabalho dinâmico e \textbf{dialógico} com a linguagem permite \textbf{que} o estudante compreenda o quanto o domínio desta é preponderante, para \textbf{que} ele tenha plenas condições de participar ativamente na sociedade.
Desse modo, tendo em vista \textbf{que} \textbf{uma} das funções da escola é garantir \textbf{que} o aluno desenvolva a habilidade da escrita, considera -se como fundamental \textbf{que} os professores de língua portuguesa estimulem constantemente o desenvolvimento dessa capacidade linguística, essencial para o processo de interação, como também para garantir \textbf{que} os estudantes sejam \textbf{autores} de seus textos/discursos.
Isso porque, conforme Geraldi ( 1993, p. 135 ) destaca, a produção de textos é considerada “ como ponto de partida de todo o processo de ensino/aprendizagem da língua [ pois ] é no texto \textbf{que} a língua se \textbf{revela} em sua \textbf{totalidade} ”.
Assim, a produção textual é \textbf{uma} das atividades \textbf{que} valoriza o papel do \textbf{sujeito} na sociedade, \textbf{uma} vez \textbf{que} é por meio de enunciados escritos \textbf{que} o indivíduo pode interagir em seu ambiente social e expor seu \textbf{posicionamento} sobre o mundo.
Nessa \textbf{direção}, assim como orientados nos documentos curriculares anteriores, ratifica -se \textbf{que} as atividades de produção de textos não devem ser desenvolvidas de forma restrita e fora do contexto, mas por meio de situações efetivas de produção de textos pertencentes a gêneros discursivos \textbf{que} circulam nos variados campos de atividade humana. \textbf{Visto} \textbf{que}, conforme o filósofo Bakhtin ( 2003 ), somente o exercício consciente das ações \textbf{que} se fazem com a linguagem e sobre ela é capaz de favorecer a instauração de \textbf{um} locutor, \textbf{que} as elabora e por elas se responsabiliza, seja nas modalidades oral, escrita e multissemiótica, ou no gênero discursivo requerido pela situação comunicativa de produção, circulação e recepção.
Diante disso, esse documento de referência curricular orienta \textbf{que} o processo de ensino e aprendizagem da escrita, em toda a Educação Básica, deve ser realizado a partir de projetos e planejamentos \textbf{que} tenham como foco a valorização de \textbf{uma} escrita autoral, constituída no processo de interação, relacionando as práticas de linguagem, tanto do texto escrito quanto do oral e do multissemiótico.
Conforme apresentado na BNCC ( 2017, p. 74 ), é possível, > [... ], por exemplo, construir \textbf{um} álbum de personagens famosas, de heróis/heroínas ou de vilões ou vilãs ; produzir \textbf{um} almanaque \textbf{que} retrate as práticas culturais da comunidade ; narrar fatos cotidianos, de forma crítica, lírica ou bem-humorada em \textbf{uma} crônica ; comentar e indicar diferentes produções culturais por meio de resenhas ou de playlists comentadas ; descrever, avaliar e recomendar ( ou não ) \textbf{um} game em \textbf{uma} resenha, gameplay ou vlog ; escrever verbetes de curiosidades científicas ; \textbf{sistematizar} dados de \textbf{um} estudo em \textbf{um} relatório ou relato multimidiático de campo ; divulgar conhecimentos específicos por meio de \textbf{um} verbete de enciclopédia digital colaborativa ; relatar fatos relevantes para a comunidade em notícias ; cobrir acontecimentos ou levantar dados relevantes para a comunidade em \textbf{uma} reportagem ; expressar posição em \textbf{uma} carta de leitor ou artigo de opinião ; denunciar situações de desrespeito aos direitos por meio de fotorreportagem, fotodenúncia, poema, lambe-lambe, microrroteiro, dentre outros.
Assim, o trabalho desenvolvido nesse viés possibilita \textbf{que} o aluno produza \textbf{uma} escrita autêntica, criativa e autônoma.
É essencial \textbf{que} os professores de língua portuguesa tenham a clareza de \textbf{que} o desenvolvimento das práticas de produção de textos abrange dimensões \textbf{inter-relacionadas} às práticas de uso e reflexão ( BNCC, 2017 ).
Antes de solicitar \textbf{que} o aluno elabore \textbf{uma} produção escrita, é necessário \textbf{que} o professor se utilize de estratégias de produção, isto é, \textbf{que} auxilie os alunos no desenvolvimento de habilidades de planejamento, revisão, edição, reescrita/redesigne e avaliação dos textos \textbf{que} produzirem, considerando a adequação desses aos contextos em \textbf{que} foram produzidos, ao modo ( escrito ou oral ; imagem estática ou em movimento etc. ), à variedade linguística e/ou semioses apropriadas a esse contexto, os enunciadores envolvidos, o gênero, o suporte, a esfera/campo de circulação, adequação à norma-padrão etc. ( BNCC, 2017 ). \textbf{Logo}, o trabalho realizado durante todo esse processo possibilitará, de fato, a \textbf{correlação} entre as práticas de uso e reflexão da linguagem. 1.1.2.4.
Conhecimentos Linguísticos e Gramaticais A linguagem é o meio \textbf{que} o ser humano utiliza para se comunicar, apresentar críticas, defender o seu ponto de vista e compartilhar conhecimentos de mundo.
Como os Parâmetros Curriculares Nacionais de Língua Portuguesa ( 1998, p. 19 ) exemplificam, > O domínio da linguagem, como atividade discursiva e cognitiva, e o domínio da língua, como sistema simbólico utilizado por \textbf{uma} comunidade linguística, são condições de possibilidade de plena participação social.
Depois de algum tempo e pouco avanço comprovado com relação ao ensino de gramática nas salas de aula, surge \textbf{um} novo documento com a intenção de padronizar parte dos conteúdos trabalhados na Educação Básica : a BNCC.
O objetivo da Base, no \textbf{que} \textbf{diz} respeito à Língua Portuguesa, é garantir aos estudantes o acesso aos saberes linguísticos necessários para a participação social e o exercício da cidadania.
A gramática assume \textbf{um} lugar muito importante nesse documento.
É possível evidenciar isso ao analisar as propostas de atividades \textbf{que} conduzem a \textbf{uma} progressão do conhecimento do aluno, \textbf{uma} vez \textbf{que} se \textbf{preocupa} com \textbf{que} domine regras e processos gramaticais, usando -os adequadamente nas diversas situações comunicativas.
É importante destacar \textbf{que} o ensino de gramática precisa acontecer para contribuir com o falante no conhecimento de sua língua materna, proporcionando -lhe \textbf{um} aprendizado harmonioso.
Vale ressaltar \textbf{que} não basta repassar regras gramaticais, já \textbf{que} esse procedimento não estimula a reflexão sobre a língua, nem contribui para a formação de leitores críticos.
Segundo Antunes ( 2007, p. 41 ), > Para ser eficaz comunicativamente não basta, portanto, saber apenas as regras específicas da Gramática, das diferentes classes de palavras, suas flexões, suas combinações possíveis, a ordem de sua colocação nas frases, seus casos de concordância, entre outras.
Tudo é necessário, mas não é suficiente.
É possível adotar \textbf{um} ensino de Língua Portuguesa diferente do tradicional e \textbf{que} aplique as sugestões propostas na BNCC.
Conforme Travaglia ( 2003, p. 17 ), a língua é \textbf{um} “ conjunto de conhecimentos linguísticos \textbf{que} o usuário tem internalizado para uso efetivo em situações concretas de interação comunicativa ”.
Dessa forma, o professor precisa desenvolver a competência comunicativa dos \textbf{aprendizes} e não apenas ensinar -lhes as regras gramaticais.
A Base mostra \textbf{que} o estudo da língua deve estabelecer relação com todos os eixos.
Sendo assim, ao utilizar a funcionalidade da gramática no texto estará utilizando a análise linguística.
Geraldi ( 1997, p. 74 ) comenta \textbf{que}, > O uso da expressão “ Análise Linguística ”, não se deve ao \textbf{mero} gosto por novas terminologias.
A análise linguística inclui tanto o trabalho sobre as questões tradicionais da gramática quanto questões amplas a propósito do texto ; adequação do texto aos objetivos pretendidos ; análise dos recursos \textbf{expressivos} utilizados ( metáforas, metonímias, paráfrases, citações, discursos direto e indireto, etc. ) ; organização e inclusão de informações etc.
Essencialmente, a prática de análise linguística não poderá limitar -se à higienização do texto do aluno em seus aspectos gramaticais e ortográficos, limitando -se a “ correções ”.
Trata -se de trabalhar com o aluno o seu texto para \textbf{que} ele atinja seus objetivos junto aos leitores a \textbf{que} se destina.
Nada deve inibir \textbf{uma} nova postura do professor no \textbf{que} \textbf{diz} respeito às práticas de análise linguística.
O tempo para estudar, preparar aulas e até mesmo oportunidade de cursos de capacitação é \textbf{imprescindível}, todavia tal necessidade precisa estar associada à \textbf{vontade} do educador de conhecer as \textbf{teorias} \textbf{que} vão lhe dar fundamentação e contribuir com sua prática, para \textbf{que} possa fazer a diferença nas aulas de Língua Portuguesa. 1.1.2.5.
Educação \textbf{Literária} Inicia -se com \textbf{uma} citação de Zilberman, por ser extrema importância \textbf{lembrar} do quanto leitura e literatura estão \textbf{imbricados} : > Da alfabetização, \textbf{tarefa} \textbf{que} a escola \textbf{desempenhou} burocraticamente desde seus inícios, passou -se à necessidade de letramento, sobretudo de letramento \textbf{literário}.
A leitura de textos apresenta -se como prática inusitada, e a literatura, em boa parte das escolas nacionais, como \textbf{um} alienígena, sobretudo nas \textbf{que} atendem os segmentos \textbf{populares}, mesmo em grandes centros urbanos ( ZILBERMAN, 2008, p. 15 ).
É comum quando se trata do uso do texto \textbf{literário} em sala de aula.
Usa -se para expor pontos de vista, para exemplificar gêneros textuais e do discurso, para estudar aspectos gramaticais da língua.
Mas, poucas vezes, o texto é usado para seu fim precípuo, principalmente no ensino fundamental : conduzir a criança de diferentes idades a \textbf{um} mundo imaginário no qual tudo “ pode ser ”.
É evidente \textbf{que} para \textbf{que} este seja bem compreendido se faz necessário o conhecimento da estrutura da língua, mas em momento algum se deve \textbf{perder} de vista as múltiplas possibilidades de leitura dessa linguagem tão particular quanto universal : a linguagem \textbf{literária}.
Isso significa necessariamente pensá -la como metáfora da própria vida, em diferentes espaços, tempos e significados.
Nesse sentido, a educação \textbf{literária} é o desvelar do texto, considerando o tripé sobre o qual se apoia o sistema \textbf{literário} : obra, \textbf{autor} e público ( CÂNDIDO, 2000 ).
Desse modo, ao pensar em sistema \textbf{literário}, na perspectiva da educação \textbf{literária}, a escola assume importância capital como rede de transmissão de ideias e de formação do leitor.
É \textbf{preciso} \textbf{lembrar} \textbf{que} na educação infantil e nos anos iniciais do fundamental a literatura aparece de forma pontual e específica na leitura deleite.
Entretanto, alguns professores ignoram a importância dessa estratégia de leitura e a subutilizam ou nem utilizam.
E quando as crianças – ainda crianças – entram nos anos finais, a ruptura com tudo o \textbf{que} significava “ escola ” e “ escolarização ” é radical, começando com o fato de \textbf{que} não há mais espaço específico para o deleite, o aprender a gostar da leitura e da literatura.
Soares ( 2004, p. 06 ) denuncia isso recorrentemente como “ precário domínio das competências de leitura e de escrita, dificultando sua inserção no mundo social e no mundo do trabalho ”.
Assim, Soares ( 2004 ), Cox ( 2004 ), Zilberman ( 2008 ), Lajolo ( 2009 ), e Fontão ( 2010 ) têm em comum o fato de considerarem \textbf{que} a escola ainda não trata adequadamente a leitura do texto \textbf{literário}, mesmo divergindo acerca do “ como ” isso pode ser feito em sala de aula para tornar o aluno competente em leitura ( literariamente alfabetizado, ou seja, letramento \textbf{literário} ).
Para a BNCC, as competências específicas para o ensino de Língua Portuguesa no ensino fundamental \textbf{dependem} em grande parte do desenvolvimento da competência leitora, especificamente \textbf{focada} em literatura na competência nove : > Envolver -se em práticas de leitura \textbf{literária} \textbf{que} possibilitem o desenvolvimento do senso estético para fruição, valorizando a literatura e outras manifestações artístico-culturais como formas de acesso às dimensões lúdicas, de imaginário e encantamento, reconhecendo o potencial transformador e humanizador da experiência com a literatura ( BNCC, 2017, p. 85 ).
Nesse sentido, Lajolo ( 2009 ) compreende \textbf{que} no texto estão presentes elementos do próprio texto e do \textbf{sujeito} leitor ( individual e coletivamente no \textbf{que} \textbf{diz} respeito ao momento \textbf{que} cada \textbf{um} lê, a história coletiva \textbf{que} \textbf{vive} e a história de cada \textbf{um} ).
Em O imaginário no poder, Held ( 1980, p. 18 ) reitera a função lúdica da literatura para crianças, \textbf{uma} vez \textbf{que}, para a autora, a leitura do “ mundo real ” passa pela leitura do “ mundo fantástico ”, presente nos contos, fábulas e outros gêneros discursivos.
É na literatura, especialmente na literatura fantástica, \textbf{que} desbloqueiam o imaginário, fazendo explodir “ estruturas fixas, estereotipadas, \textbf{que} transformam o universo cotidiano, \textbf{que} criam \textbf{um} passado, \textbf{um} presente, \textbf{um} futuro e \textbf{uma} dinâmica criativa irreversível ”.
Segundo Fontão ( 2010, p. 4 ) : > Se pensarmos a escola do passado e a escola \textbf{contemporânea}, vamos perceber \textbf{que} o desenvolvimento de \textbf{um} leitor proficiente não se faz somente com \textbf{um} tipo de texto e sim com \textbf{um} conjunto de livros e textos \textbf{que} apresentam linguagens variadas, como também, variados devem ser os \textbf{métodos} de desenvolvimento de leitura. [ Assim ] não basta ao aluno ler \textbf{um} clássico ou mesmo \textbf{um} livro didático \textbf{que} traga fragmentos de textos de obras-primas da literatura universal e brasileira, mas sim ler pelo gosto do ler, pela aventura do ler, pelo desejo de penetrar no mundo da leitura daquele determinado texto, daquela determinada história, apreendendo sobre a estética, as regras, características e parâmetros, a fim de descortinar e sanar o nível de dificuldade em realizar a leitura de diferentes gêneros textuais, produzir e interpretar textos.
Essa forma de \textbf{ver} a literatura no ensino fundamental, presente também na BNCC, como já foi \textbf{visto}, faz pensar em \textbf{um} problema/solução : o professor deve ser leitor competente e amante da literatura para \textbf{que} o estudante também o seja.
Pound ( 2006, p. 80 ) acredita \textbf{que} “ o professor \textbf{ideal} seria o \textbf{que} examinasse qualquer obra-prima \textbf{que} estivesse apresentando a seus alunos quase como se nunca a tivesse \textbf{visto} antes ” ( \textbf{grifo} do \textbf{autor} ). \textbf{Talvez}, assim, parte do desafio seja instrumentalizar os professores para desvelá -lo da literatura infanto-juvenil e pensá -la como \textbf{uma} grande aventura em \textbf{que}, juntos, professor e aluno, descubram o mundo.

% matched lemmas: abordagem, aprendiz, ator, atual, autor, basear, conseguir, contemporâneo, convir, correlação, depender, desempenhar, dialógico, direção, dizer, exigir, expressivo, focar, grifo, historicamente, ideal, imbricar, imprescindível, inter-relacionar, interessante, lembrar, literário, logo, mero, método, perder, pontuar, popular, posicionamento, postar, preciso, preocupar, que, revelar, sistematizar, sujeito, século, talvez, tarefa, teoria, totalidade, tópico, um, ver, viver, vontade
\end{textsample}
